\documentclass{article}
\usepackage{arxiv}
% For french
\usepackage[utf8]{inputenc}
\usepackage[T1]{fontenc}
\usepackage{amsthm,amsmath,natbib}
\usepackage{amsfonts,amssymb}
\RequirePackage[colorlinks,citecolor=blue,urlcolor=blue]{hyperref}
\usepackage[toc,page]{appendix}

% For tikz pictures
\usepackage{tikz}
\usetikzlibrary{matrix,chains,positioning,decorations.pathreplacing,arrows}
\usepackage{neuralnetwork}
\usepackage{pgfplots}
\pgfplotsset{compat=1.15}
\usepackage{algorithm}
\usepackage[]{algorithmicx}
\usepackage{algpseudocode}
\usepackage{afterpage}
\usepackage{bbm}
\usepackage{setspace}
\usepackage{comment}
\usepackage[normalem]{ulem}
%\doublespacing

\newcommand\independent{\protect\mathpalette{\protect\independenT}{\perp}}
\def\independenT#1#2{\mathrel{\rlap{$#1#2$}\mkern2mu{#1#2}}}

% Additional definitions
\def\qadj{\hat{q}_{\mathrm{adj}}}
\def\tt{\mathbf t}
\def\z{\mathbf z}
\def\X{\mathbf X}
\def\inv{^{-1}}
\def\H{\mathbf H}
\def\x{\mathbf x}
\def\xs{\underline {\mathbf x}}
\def\Xs{\underline { \mathbf X} }
\def\r{\mathbf r}
\def\s{\mathbf s}
\def\y{\mathbf y}
\def\u{\mathbf u}
\def\v{\mathbf v}
\def\lmax{\lambda_\mathrm{max}}
\def\emax{\mathbf{e}_\mathrm{max}}
\def\emin{\mathbf{e}_\mathrm{min}}
\def\lmin{\lambda_\mathrm{min}}
\def\bbeta{\boldsymbol \beta}
\def\ggamma{\boldsymbol \gamma}
\def\intinf{\int_\infty^\infty}
\def\sumi{\sum_{i=1}^n}
\def\t{^\top}
\def\eps{\boldsymbol \varepsilon}
\def\KL{\mathbb {KL}}
\def\E{\mathbb E}
\def\V{\mathbb V}
\def\N{\mathcal N}
\def \real{\rm I\!R}
\def \I{\mathbf I}
\def \L {\mathcal{L}}

% define alphabetical enumeration in \bgin{enumerate}
\renewcommand{\theenumi}{\alph{enumi}}
\renewcommand{\labelenumi}{\theenumi)}
\newcommand{\norm}[1]{\left\lVert#1\right\rVert}

\usepackage{mathtools}
\DeclarePairedDelimiter{\ceil}{\lceil}{\rceil}
\DeclareMathOperator*{\argmax}{argmax~} % for subscripts and superscripts
\DeclareMathOperator*{\argmin}{argmin~} % for subscripts and superscripts
\DeclareMathOperator{\sign}{sign} % for functions


\newtheorem{theorem}{Theorem}
\newtheorem{acknowledgement}[theorem]{Acknowledgement}
%\newtheorem{algorithm}[theorem]{Algorithm}
\newtheorem{axiom}[theorem]{Axiom}
\newtheorem{case}[theorem]{Case}
\newtheorem{claim}[theorem]{Claim}
\newtheorem{conclusion}[theorem]{Conclusion}
\newtheorem{condition}[theorem]{Condition}
\newtheorem{conjecture}[theorem]{Conjecture}
\newtheorem{corollary}[theorem]{Corollary}
\newtheorem{criterion}[theorem]{Criterion}
\newtheorem{definition}[theorem]{Definition}
\newtheorem{example}[theorem]{Example}
\newtheorem{exercise}[theorem]{Exercise}
\newtheorem{lemma}[theorem]{Lemma}
\newtheorem{notation}[theorem]{Notation}
\newtheorem{problem}[theorem]{Problem}
\newtheorem{proposition}[theorem]{Proposition}
\newtheorem{remark}[theorem]{Remark}
\newtheorem{solution}[theorem]{Solution}
\newtheorem{summary}[theorem]{Summary}
%\newenvironment{proof}[1][Proof]{\noindent\textbf{#1.} }{\ \rule{0.5em}{0.5em}}


\title{Generating Synthetic Data for Treatment Effect Estimation}

\author{
  Mouloud Belbahri \\
  mouloud.belbahri@tdassurance.com \\
  Department of Analytics and Modeling\\
  Advanced Projects
}


\begin{document}
\maketitle

\begin{abstract}
This short report develops a method to evaluate different treatment estimation models using synthetic data generated in a similar way to the simulation studies of \cite{tian2014simple}.
\end{abstract}

\keywords{causal inference \and simulations}


%\section{Data generation}

Synthetic data generation depends on the following four elements.

\begin{enumerate}
    \item The number $n$ of observations in the sample, and the number $p$ of predictors observed for each observation.
    \item The distribution of the predictors. We draw predictors from a multivariate Gaussian distribution.
    \item The mean effect function $\mu(\X)$ and the treatment effect function $\tau(\X)$. We vary the size of main effects relative to treatment effects.
    \item The noise level $\sigma^2$. This corresponds to the conditional variance of the outcome given $\X$ and $T$. 
\end{enumerate}

Given the elements above, our data generation model is, for $i=1,...,n$ :

\begin{align*}
    &\X_i \sim \mathcal{N}_p (0, \Sigma) \\
    &T_i \sim \mathrm{Bernoulli}(\theta)\\
    &Y_i^* \mid \X_i, T_i \sim \mathcal{N}(\mu(\X_i) + T_i\tau(\X_i), \sigma^2).
\end{align*}

Define $Y_i$ to be the binary outcome as

\begin{align*}
    Y_i \mid \X_i, T_i = \mathbbm{1}(Y_i^* > 0 \mid \X_i, T_i).
\end{align*}

Hence, $Y_i$ follows a Bernoulli distribution with parameter $\mathrm{Pr}(Y_i^* > 0 \mid \X_i, T_i)$ and it is easy to recover the ``true'' treatment effect $u_i^*$ as follows

\begin{align*}
    u_i^* &= \mathrm{Pr}(Y_i^* > 0 \mid \X_i, T_i = 1) - \mathrm{Pr}(Y_i^* > 0 \mid \X_i, T_i = 0)\\
    u_i^* &= \Phi \Bigl( \frac{\mu(\X_i) + \tau(\X_i)}{\sigma} \Bigr) - \Phi \Big( \frac{\mu(\X_i)}{\sigma} \Big) 
\end{align*}
where $\Phi(.)$ is the standard Gaussian cumulative distribution function. Note that $\mu(.)$ is a nuisance function and $\tau(.)$ is the function that controls the treatment effect. In fact, the sign of $\tau(.)$ and the sign of the treatment effect for a given observation coincide.

We set $\theta$ to $0.5$, which is the case in our real data. For observational data settings, one can define $\theta$ as a function of a subset of the covariates, e.g. $\theta= 1 / \{1 + e^{-(X_1 + X_2)} \}$. 

The variance-covariance matrix elements $\Sigma_{ij}$ are defined as follows:

\[\Sigma_{ij} = \begin{cases} 
      \rho & i\neq j \\
      1 & i=j
   \end{cases}.
\]

Finally, the mean and treatment effects functions are defined as follows:

$$\mu(\X) = \sum_{j=1}^J \beta_j X_j + \sum_{j=1}^{p-1} \omega_j X_j X_{j+1},$$ 


$$\tau(\X) = \sum_{j=1}^K \gamma_j X_j$$

where $J\leq p$ and $K \leq p$. We set $\beta_j = \beta (-1)^{j}$, $\omega_j = \omega (-1)^j$ and we draw $\gamma_j$ independently from $\mathcal{N}(0,\gamma)$. The constants $\beta$, $\omega$ and $\gamma$ are hyper-parameters.





\bibliographystyle{plainnat}
\bibliography{mybib}



\end{document}